\documentclass[10pt,a4paper]{amsart}
\usepackage{amsmath,amssymb,amsthm,graphicx}
% Retained preamble for the Stacks Project chapter {CH} (shared project preamble).
% The chapter file pinned-original-source.tex only carries \input{{preamble}} on its first line,
% so no preamble text lives inside the chapter itself.  The Stacks Project uses one single shared
% preamble for every chapter: the snapshot stages it three times (sources/stacks-chow/preamble.tex,
% sources/stacks-etale-cohomology/preamble.tex and focused-seeds/AGP/sources/stacks-snapshot-20261011/
% preamble.tex) and those copies are identical in size.  The block below is copied mechanically,
% line for line, from sources/stacks-chow/preamble.tex lines 14-18 and 154-217: the Xy-pic setup and
% the native theorem declarations together with the mathematical macros.  Nothing was added, removed
% or redefined; the wrapper class and the amsmath/amssymb/amsthm/graphicx base packages are supplied
% by the assembler.
% For commutative diagrams we use Xy-pic
% Chapter: stacks-spaces-more-morphisms.tex -- More on Morphisms of Spaces.
\usepackage[all]{xy}

% We use 2cell for 2-commutative diagrams.
\xyoption{2cell}

% Theorem environments.
%
\theoremstyle{plain}
\newtheorem{theorem}[subsection]{Theorem}
\newtheorem{proposition}[subsection]{Proposition}
\newtheorem{lemma}[subsection]{Lemma}

\theoremstyle{definition}
\newtheorem{definition}[subsection]{Definition}
\newtheorem{example}[subsection]{Example}
\newtheorem{exercise}[subsection]{Exercise}
\newtheorem{situation}[subsection]{Situation}

\theoremstyle{remark}
\newtheorem{remark}[subsection]{Remark}
\newtheorem{remarks}[subsection]{Remarks}

\numberwithin{equation}{subsection}

% Macros
%
\def\lim{\mathop{\mathrm{lim}}\nolimits}
\def\colim{\mathop{\mathrm{colim}}\nolimits}
\def\Spec{\mathop{\mathrm{Spec}}}
\def\Hom{\mathop{\mathrm{Hom}}\nolimits}
\def\Ext{\mathop{\mathrm{Ext}}\nolimits}
\def\SheafHom{\mathop{\mathcal{H}\!\mathit{om}}\nolimits}
\def\SheafExt{\mathop{\mathcal{E}\!\mathit{xt}}\nolimits}
\def\Sch{\mathit{Sch}}
\def\Mor{\mathop{\mathrm{Mor}}\nolimits}
\def\Ob{\mathop{\mathrm{Ob}}\nolimits}
\def\Sh{\mathop{\mathit{Sh}}\nolimits}
\def\NL{\mathop{N\!L}\nolimits}
\def\CH{\mathop{\mathrm{CH}}\nolimits}
\def\proetale{{pro\text{-}\acute{e}tale}}
\def\etale{{\acute{e}tale}}
\def\QCoh{\mathit{QCoh}}
\def\Ker{\mathop{\mathrm{Ker}}}
\def\Im{\mathop{\mathrm{Im}}}
\def\Coker{\mathop{\mathrm{Coker}}}
\def\Coim{\mathop{\mathrm{Coim}}}

% Boxtimes
%
\DeclareMathSymbol{\boxtimes}{\mathbin}{AMSa}{"02}

%
% Macros for moduli stacks/spaces
%
\def\QCohstack{\mathcal{QC}\!\mathit{oh}}
\def\Cohstack{\mathcal{C}\!\mathit{oh}}
\def\Spacesstack{\mathcal{S}\!\mathit{paces}}
\def\Quotfunctor{\mathrm{Quot}}
\def\Hilbfunctor{\mathrm{Hilb}}
\def\Curvesstack{\mathcal{C}\!\mathit{urves}}
\def\Polarizedstack{\mathcal{P}\!\mathit{olarized}}
\def\Complexesstack{\mathcal{C}\!\mathit{omplexes}}
% \Pic is the operator that assigns to X its picard group, usage \Pic(X)
% \Picardstack_{X/B} denotes the Picard stack of X over B
% \Picardfunctor_{X/B} denotes the Picard functor of X over B
\def\Pic{\mathop{\mathrm{Pic}}\nolimits}
\def\Picardstack{\mathcal{P}\!\mathit{ic}}
\def\Picardfunctor{\mathrm{Pic}}
\def\Deformationcategory{\mathcal{D}\!\mathit{ef}}
\usepackage{cleveref}
\crefname{originalreference}{original reference}{original references}
\Crefname{originalreference}{Original reference}{Original references}
\makeatletter
\newlabel{descent-lemma-pullback-property-local-target}{{04QU}{source}{}{sourcepointer-1}{}}
\newlabel{descent-lemma-pullback-property-local-target@cref}{{[originalreference][0][]04QU}{[1][0][]source}{}{}{}}
\makeatother
\title{Testing a property of a representable morphism of algebraic spaces after an etale cover}
\author{The Stacks Project Authors}
\date{stacks-project:spaces-more-morphisms.tex:local-snapshot-20261011; GNU FDL 1.2 or later}
\begin{document}
\maketitle
\noindent\textit{Editorial note.} Counted span is the author's own proof environment, source lines 10477-10520, immediately following the statement at 10453-10475.\par\medskip
\noindent\textit{Editorial note.} One cross-chapter pointer (Descent) is recorded in assembly.external\_references.\par\medskip
\noindent\textit{Editorial note.} The statement quantifies over a property P fixed by the reader; that wording is retained verbatim in the statement span.\par\medskip
\noindent\textit{Allowed standard prerequisites (editorial).}\begin{itemize}
\item Properties of morphisms etale local on the target. the property P fixed in the statement, with Descent, lemma descent-lemma-pullback-property-local-target cited at proof lines 10508-10509. Supplies that P is preserved under etale base change and can be tested after an etale covering of the target.
\item Representability of a morphism of algebraic spaces. used throughout: representable morphisms are those whose base change to a scheme is a scheme morphism. Lets the diagrams with V a scheme and V -> Y etale be used as the testing devices.
\item Surjectivity and etaleness of the diagonal cover. the morphism U x\_X U' -> U' appearing at proof lines 10511-10514. Provides the etale covering of U' which finishes the implication.
\end{itemize}\medskip
\ifdefined\hypertarget\hypertarget{sourcepointer-1}{}\fi
\noindent\textit{Source reference.} Original reference 04QU: focused-seeds/geometry/S9-spaces-representable-etale-local-target/ref-tags.txt, source lines 4--4. Navigation pointer only: the counted proof cites this statement by \textbackslash{}ref as a standing input; the deduction the counted proof has to make is carried out inside the counted span, and no part of that deduction is delegated to this pointer.\par\medskip
\begin{lemma}
\label{lemma-representable-etale-local-target}
Let $\mathcal{P}$ be a property of morphisms of schemes which is \'etale
local on the target. Let $S$ be a scheme.
Let $f : X \to Y$ be a representable morphism of algebraic spaces over $S$.
Consider commutative diagrams
$$
\xymatrix{
X \times_Y V \ar[d] \ar[r] & V \ar[d] \\
X \ar[r]^f & Y
}
$$
where $V$ is a scheme and $V \to Y$ is \'etale.
The following are equivalent
\begin{enumerate}
\item for any diagram as above the projection $X \times_Y V \to V$
has property $\mathcal{P}$, and
\item for some diagram as above with $V \to Y$ surjective
the projection $X \times_Y V \to V$ has property $\mathcal{P}$.
\end{enumerate}
If $X$ and $Y$ are representable, then this is also equivalent to
$f$ (as a morphism of schemes) having property $\mathcal{P}$.
\end{lemma}

\begin{proof}
Let us prove the equivalence of (1) and (2).
The implication (1) $\Rightarrow$ (2) is immediate.
Assume
$$
\xymatrix{
X \times_Y V \ar[d] \ar[r] & V \ar[d] \\
X \ar[r]^f & Y
}
\quad\quad
\xymatrix{
X \times_Y V' \ar[d] \ar[r] & V' \ar[d] \\
X \ar[r]^f & Y
}
$$
are two diagrams as in the lemma. Assume $V \to Y$ is
surjective and $X \times_Y V \to V$ has property $\mathcal{P}$.
To show that (2) implies (1) we have to prove that
$X \times_Y V' \to V'$ has $\mathcal{P}$. To do
this consider the diagram
$$
\xymatrix{
X \times_Y V \ar[d] &
(X \times_Y V) \times_X (X \times_Y V') \ar[l] \ar[d] \ar[r] &
X \times_Y V' \ar[d] \\
V &
V \times_Y V' \ar[l] \ar[r] &
V'
}
$$
By our assumption that $\mathcal{P}$ is \'etale local on the source,
we see that $\mathcal{P}$ is preserved under \'etale base change, see
Descent, Lemma \ref{descent-lemma-pullback-property-local-target}.
Hence if the left vertical arrow has $\mathcal{P}$ the so does
the middle vertical arrow. Since $U \times_X U' \to U'$ is surjective
and \'etale (hence defines an \'etale covering of $U'$)
this implies (as $\mathcal{P}$ is assumed local for the \'etale topology
on the target) that the left vertical arrow has $\mathcal{P}$.

\medskip\noindent
If $X$ and $Y$ are representable, then we can take
$\text{id}_Y : Y \to Y$ as our \'etale covering to see the
final statement of the lemma is true.
\end{proof}

\end{document}
